\section{A series for the hexadecimal digits of $\pi$}

A factor $16^k$ in the denominator is a shift in base $16$. The series below isolates that shift.

\begin{theorem}[Bailey, Borwein and Plouffe, 1995]\label{thm:bbp}
For $j\ge1$, let $S_j=\sum_{k\ge0}\frac1{16^k(8k+j)}$. Then $\pi=4S_1-2S_4-S_5-S_6$.
\end{theorem}

\begin{proof}
For $k\ge0$, $\int_0^{1/\sqrt2}x^{8k+j-1}\dd x=\frac{2^{-j/2}}{16^k(8k+j)}$. Summing over $k$ (the terms are positive, so the sum and the integral can be exchanged),
\[
S_j=2^{j/2}\int_0^{1/\sqrt2}\frac{x^{j-1}}{1-x^8}\dd x.
\]
So
\[
4S_1-2S_4-S_5-S_6=\int_0^{1/\sqrt2}\frac{4\sqrt2-8x^3-4\sqrt2x^4-8x^5}{1-x^8}\dd x,
\]
and the substitution $x=y/\sqrt2$ turns this into
\[
\int_0^1\frac{16(4-2y^3-y^4-y^5)}{16-y^8}\dd y=\int_0^1\frac{16(y-1)}{(y^2-2)(y^2-2y+2)}\dd y,
\]
since $16-y^8=(4-y^4)(4+y^4)$, $4-y^4=-(y^2-2)(y^2+2)$, $4+y^4=(y^2-2y+2)(y^2+2y+2)$, and $4-2y^3-y^4-y^5=(1-y)(y^2+2)(y^2+2y+2)$, as one checks by expanding. In partial fractions,
\[
\frac{16(y-1)}{(y^2-2)(y^2-2y+2)}=\frac{4y}{y^2-2}-\frac{4(y-1)}{(y-1)^2+1}+\frac4{(y-1)^2+1},
\]
(multiply out: the numerator is $4y(y^2-2y+2)-(4y-8)(y^2-2)=16y-16$). The three integrals over $[0,1]$ are $2\ln\frac12$, $-2\ln\frac12$ and $4\bigl(\arctan0-\arctan(-1)\bigr)=\pi$. Their sum is $\pi$.
\end{proof}

\begin{proposition}[digit extraction]\label{prop:digits}
For $d\ge0$ and $j\ge1$, let
\[
T_j(d)=\sum_{k=0}^d\frac{16^{d-k}\bmod(8k+j)}{8k+j}+\sum_{k>d}\frac{16^{d-k}}{8k+j}.
\]
Then the fractional part of $16^d\pi$ equals the fractional part of $4T_1(d)-2T_4(d)-T_5(d)-T_6(d)$. In particular the hexadecimal digits of $\pi$ from position $d+1$ on can be
computed with $O(d\log d)$ arithmetic operations on numbers of size $O(\log d)$ bits, without computing the previous digits.
\end{proposition}

\begin{proof}
For $k\le d$, $\frac{16^{d-k}}{8k+j}$ differs from $\frac{16^{d-k}\bmod(8k+j)}{8k+j}$ by an integer. So $16^dS_j-T_j(d)$ is an integer, and as the coefficients $4,-2,-1,-1$ are integers,
$16^d\pi-(4T_1-2T_4-T_5-T_6)$ is an integer by Theorem~\ref{thm:bbp}. Each residue $16^{d-k}\bmod(8k+j)$ is computed by repeated squaring modulo $8k+j$ with $O(\log d)$ multiplications of
numbers below $8d+8$, and there are $O(d)$ of them; the tail is less than $\sum_{m\ge1}16^{-m}=\frac1{15}$ and a few terms give it to any fixed accuracy.
\end{proof}

In floating point, each $T_j(d)$ is a sum of $d$ numbers in $[0,1)$, computed with relative error about $10^{-16}$; the error in the final fractional part
is then about $d\cdot10^{-16}$, so for $d=10^6$ one gets about eight correct hexadecimal digits. Direct computation gives the digits of $\pi$ starting at position $10^6$ in a few seconds:
\texttt{26C65E52}\dots, the same digits that Bailey, Borwein and Plouffe printed in 1997.



\section{The Borwein integrals}

In 2001 David and Jonathan Borwein published a series of integrals that has fooled many computer algebra systems and their users. Write
$\operatorname{sinc}x=\frac{\sin x}x$ (with $\operatorname{sinc}0=1$). Then
\[
\int_0^\infty\operatorname{sinc}x\dd x=\frac\pi2,\qquad\int_0^\infty\operatorname{sinc}x\operatorname{sinc}\frac x3\dd x=\frac\pi2,\qquad
\int_0^\infty\operatorname{sinc}x\operatorname{sinc}\frac x3\operatorname{sinc}\frac x5\dd x=\frac\pi2,
\]
and so on up to the factor $\operatorname{sinc}\frac x{13}$: seven integrals, all equal to $\frac\pi2$. The eighth,
\[
\int_0^\infty\operatorname{sinc}x\operatorname{sinc}\frac x3\cdots\operatorname{sinc}\frac x{15}\dd x
=\frac{467\,807\,924\,713\,440\,738\,696\,537\,864\,469}{935\,615\,849\,440\,640\,907\,310\,521\,750\,000}\,\pi,
\]
is smaller than $\frac\pi2$ by about $2.31\times10^{-11}$. The explanation is a piece of probability.

\begin{theorem}\label{conjectures:thm:borwein}
Let $a_1,\dots,a_n>0$, let $U_1,\dots,U_n$ be independent and uniform on $[-1,1]$, and let $S=a_1U_1+\dots+a_nU_n$. Then
\[
\int_0^\infty\operatorname{sinc}x\prod_{k=1}^n\operatorname{sinc}(a_kx)\dd x=\frac\pi2\,\mathbb P\bigl(|S|<1\bigr).
\]
In particular the integral equals $\frac\pi2$ if $a_1+\dots+a_n\le1$, and is strictly smaller if $a_1+\dots+a_n>1$.
\end{theorem}
\begin{proof}
For a uniform $U$ on $[-1,1]$ and real $t$, $\mathbb E[\cos(tU)]=\frac12\int_{-1}^1\cos(tu)\dd u=\operatorname{sinc}t$. Since $S$ is symmetric,
$\mathbb E[\sin(xS)]=0$, and by independence
\[
\mathbb E[\cos(xS)]=\operatorname{Re}\mathbb E[e^{ixS}]=\operatorname{Re}\prod_k\mathbb E[e^{ixa_kU_k}]=\prod_k\operatorname{sinc}(a_kx).
\]
For $T>0$, Fubini's theorem on $[0,T]\times\Omega$ (the integrand is bounded by $1$) gives
\[
\int_0^T\operatorname{sinc}x\prod_k\operatorname{sinc}(a_kx)\dd x=\mathbb E\int_0^T\frac{\sin x\cos(xS)}x\dd x=\frac12\,\mathbb E\bigl[D_T(1+S)+D_T(1-S)\bigr],
\]
where $D_T(c)=\int_0^T\frac{\sin(cx)}x\dd x$, by $2\sin x\cos(xS)=\sin((1+S)x)+\sin((1-S)x)$. Substituting $u=|c|x$, $D_T(c)=\operatorname{sgn}(c)\int_0^{|c|T}\frac{\sin u}u\dd u$.
The function $t\mapsto\int_0^t\frac{\sin u}u\dd u$ is continuous on $[0,\infty[$ and tends to $\frac\pi2$ (Theorem~\ref{thm:dirichletint}), so it is bounded, say by $B$.
Hence $|D_T(c)|\le B$ for all $c$ and $T$, and $D_T(c)\to\frac\pi2\operatorname{sgn}(c)$ as $T\to\infty$. By dominated convergence,
\[
\int_0^\infty\operatorname{sinc}x\prod_k\operatorname{sinc}(a_kx)\dd x=\frac\pi4\,\mathbb E\bigl[\operatorname{sgn}(1+S)+\operatorname{sgn}(1-S)\bigr].
\]
The bracket equals $2$ if $|S|<1$, $1$ if $|S|=1$ and $0$ if $|S|>1$. Since $S$ has a density when $n\ge1$, $\mathbb P(|S|=1)=0$, and the value is
$\frac\pi2\mathbb P(|S|<1)$. Finally $|S|\le a_1+\dots+a_n$, so $\mathbb P(|S|<1)=1$ when the sum is at most $1$. When the sum exceeds $1$, the event
$S>1$ contains a small box of positive volume around $(1,\dots,1)$, so $\mathbb P(|S|<1)<1$.
\end{proof}

For the Borwein integrals, $a_k=\frac1{2k+1}$. The sum $\frac13+\frac15+\dots+\frac1{13}=\frac{43\,024}{45\,045}=0.955\ldots$ is less than $1$, but adding $\frac1{15}$
gives $1.0218\ldots$ (Figure~\ref{conjectures:fig:borwein}). The pattern breaks exactly when the sum crosses $1$, and the deficit is tiny because $S$ can exceed $1$ only
when all seven $U_k$ are very close to $\pm1$ with the same sign.

\begin{corollary}\label{conjectures:cor:borweindeficit}
Let $\varepsilon=\sum_{k=1}^7\frac1{2k+1}-1=\frac{982}{45\,045}$. Then
\[
\frac\pi2-\int_0^\infty\prod_{k=0}^7\operatorname{sinc}\frac x{2k+1}\dd x=\pi\,\frac{3\cdot5\cdot7\cdots15}{7!\,2^7}\,\varepsilon^7\approx2.31\times10^{-11}.
\]
\end{corollary}
\begin{proof}
By Theorem~\ref{conjectures:thm:borwein} and symmetry, the deficit is $\frac\pi2\mathbb P(|S|\ge1)=\pi\,\mathbb P(S\ge1)$. Put $V_k=1-U_k$, independent and uniform on $[0,2]$. Then
$S=\sum a_k-\sum a_kV_k$, and $S\ge1$ means $\sum a_kV_k\le\varepsilon$. This forces $V_k\le\frac\varepsilon{a_k}\le15\varepsilon<2$ for every $k$, so the event is
the simplex $\{v_k\ge0,\ \sum a_kv_k\le\varepsilon\}$, which lies inside $[0,2]^7$. Its volume is $\frac{\varepsilon^7}{7!\,a_1\cdots a_7}$, and the density of
$(V_1,\dots,V_7)$ is $2^{-7}$. Multiplying, $\mathbb P(S\ge1)=\frac{\varepsilon^7}{7!\,2^7a_1\cdots a_7}$, which is the formula. Computed in exact arithmetic, it
agrees with the Borweins' fraction above.
\end{proof}

\begin{figure}[ht]
\centering
\fig{conjectures-ch01-borwein}
\caption{Left: the densities of $S_n=\frac{U_1}3+\dots+\frac{U_n}{2n+1}$ for $n=1,2,3,6$, computed by convolution; all live inside $[-1,1]$, so the first seven
Borwein integrals equal $\frac\pi2$. Right: the partial sums $\frac13+\dots+\frac1{2n+1}$. The seventh, $1.0218$, is the first to exceed $1$: from then on $S_n$
can leave $[-1,1]$, with probability $7.4\times10^{-12}$ on each side for $n=7$.}
\label{conjectures:fig:borwein}
\end{figure}
